%=============================================================
%  HUM-470 – Aquaculture and Marine Biodiversity pressures
%  Formatted to match the EPFL HUM-470 report style
%=============================================================
\documentclass[12pt, a4paper]{report}

% --- Packages --------------------------------
\usepackage[T1]{fontenc}
\usepackage[utf8]{inputenc}
\usepackage{lmodern}
\usepackage[english]{babel}
\usepackage{geometry}
\geometry{top=2.5cm, bottom=2.5cm, left=2.5cm, right=2.5cm}

\usepackage{graphicx}
\usepackage{tabularx} 
\usepackage{booktabs}
\usepackage{array}
\usepackage{amsmath}
\usepackage{amssymb}
\usepackage{url}
\usepackage[hidelinks]{hyperref}
\usepackage{hyperref}
\hypersetup{
    colorlinks=true,
    linkcolor=black,
    citecolor=black,
    urlcolor=blue
}
\usepackage{xcolor}
\usepackage{titlesec}
\usepackage{setspace}
\usepackage{parskip}
\usepackage[authoryear]{natbib}
\usepackage{caption}
\usepackage{float}
\usepackage{microtype}
\usepackage{enumitem}
\usepackage{fancyhdr}
\usepackage{emptypage}
\usepackage{etoolbox}
\patchcmd{\chapter}{\clearpage}{}{}{}
\usepackage{titlesec}
\titleformat{\chapter}[block]{\normalfont\huge\bfseries}{\chaptertitlename\ \thechapter}{0pt}{\huge}
\titlespacing*{\chapter}{0pt}{10pt}{0pt}

% --- Header / Footer ---------------------------
\pagestyle{fancy}
\fancyhf{}
\fancyhead[L]{\small HUM-470 Economic growth and sustainability II}
\fancyhead[R]{\small \today}
\fancyfoot[C]{\thepage}
\renewcommand{\headrulewidth}{0.4pt}

% --- Chapter title style ------------------------
\titleformat{\chapter}[display]
  {\normalfont\huge\bfseries}
  {\chaptertitlename\ \thechapter}{20pt}{\Huge}

% --- Glossary helper ---------------------------
\newcommand{\glossitem}[2]{\noindent\textbf{#1} \quad #2 \par\smallskip}

%=============================================================
\begin{document}

% --- Title page ------------------------------
\begin{titlepage}
  \thispagestyle{empty}

  % Place epfl-logo.png in the same folder, or comment out this line
  \noindent
  \includegraphics[width=7cm]{epfl-logo.png}

  \vspace{1em}
  \noindent\rule{\textwidth}{0.4pt}\\[0.5em]
  \noindent\textbf{HUM-470 Economic growth and sustainability II}
  \hfill
  \textbf{Spring 2025 - 2026}
  \\[-0.5em]
  \noindent\rule{\textwidth}{0.4pt}

  \vspace{3cm}
  \begin{center}
    {\LARGE\bfseries Aquaculture and Marine Biodiversity\\[0.4em]
     pressures in Norway's Salmon Industry}
  \end{center}

  \vspace{1cm}
  \begin{center}
    \textit{Final Report, Group 02}
  \end{center}

  \vfill

  \begin{center}
    \textbf{Authors}\\[0.4em]
    Victor Legrand\\
    Thomas Hetier\\
    Elhadji Ousmane Ndiaye\\
    Dominic Bazina-Grolinger
  \end{center}

  \vspace{2em}
  \noindent\rule{\textwidth}{0.4pt}\\[0.3em]
  \noindent
  \begin{minipage}[t]{0.5\textwidth}
    \textbf{Supervisors:}\\
    Dr.\ Sascha Nick
  \end{minipage}%
  \begin{minipage}[t]{0.5\textwidth}
    \raggedleft
    \textbf{Supervising Teacher:}\\
    Prof.\ Philippe Thalmann
  \end{minipage}
  \\[0.3em]
  \noindent\rule{\textwidth}{0.4pt}
\end{titlepage}

% --- Table of Contents --------------------------
\tableofcontents
\newpage


\setcounter{page}{1}

\setcounter{chapter}{0}
\chapter{Introduction}

Norway is home to one of the most productive and economically important
aquaculture industries in the world. Atlantic salmon farming has grown
dramatically along the Norwegian coastline over the past five decades,
transforming from a small-scale coastal activity into a global export industry worth billions of Norwegian kroner annually. Today, Norway produces
roughly half of the world's farmed Atlantic salmon, with production concentrated in sea-based net pen farms located in the fjords and coastal waters stretching from the south to the far north of the country.

This economic success also brings environmental costs. Firstly, salmon are carnivorous fish and they require food that is rich in protein and omega-3 fatty acids. Historically, a large share of this protein has come from wild-caught fish such as anchovies, herring, and sand eels. Harvesting these species at large scale to feed farmed salmon puts pressure on the marine ecosystem by reducing food available for seabirds, marine mammals, larger fish, and other natural predators. Secondly, the high densities of salmon in net pens create ideal conditions for sea lice to spread. Sea lice are small parasitic crustaceans that attach to fish and feed on their skin. They also have the potential to spread to and weaken wild Atlantic salmon in nearby rivers and fjords. Controlling sea lice has historically required the use of chemical treatments, many of which are toxic to crustaceans such as lobsters and shrimp when released into
the surrounding water.

Both of these pressures, use of wild fish for feed and chemical pollution from lice treatments are closely tied to historical salmon farming operations and the large scale of the industry. As production has expanded, an important question is whether these impacts have also increased, or if improvements in technology and regulation have helped limit them, making the industry more sustainable.


\section{Research Questions}

Our analysis into this industry is structured around two research questions: 

\begin{enumerate}
    \item \textbf{Has absolute decoupling occurred?} Have the two biodiversity pressures we consider: total wild fish extraction \(W\) and total persistent antiparasitic discharge \(A_{out}\) declined in absolute terms over the study period, even though the industry has experienced massive growth in recent years?

    \item \textbf{What factors explain observed trends?} Using IPAT
    decompositions, what specific drivers: feed efficiency, ingredient
    substitution, treatment type, lice management are responsible for the
    observed changes in \(W\) and \(A_{out}\)?
\end{enumerate}

\section{Structure of the Report} 

The report is organized as follows. Section~2 provides the scientific
and industrial context, covering the growth of Norwegian salmon
aquaculture, the three main impact mechanisms on marine biodiversity,
and a brief review of existing studies on decoupling in this sector.
Section~3 presents the theoretical foundations of the IPAT approach
and the structure of both decompositions. Section~4 describes our
data sources and the key assumptions we make. Sections 5 and 6 present and discuss the results of IPAT 1 (wild fish extraction) and IPAT 2 (chemical discharge) respectively. Section 7 then introduces the Norwegian regulatory framework, using the results of both decompositions as a basis for explaining which policies and incentives shaped the observed outcomes. Section 8 offers a broader discussion, comparing the two decouplings and evaluating the limits of our models. Section~9 draws recommendations for policy and industry practice, and Section~10
concludes.




%==================================================
% 2. SCIENTIFIC AND INDUSTRIAL CONTEXT
%==================================================
\chapter{Scientific and Industrial Context}

\section{Norway's Salmon Aquaculture: Growth and Economic Significance}

Atlantic salmon farming in Norway has been rapidly growing and changing over the last few decades. What began in the 1970s as a small-scale activity along the Norwegian coast has grown into one of the most productive and economically significant aquaculture industries in the world \citep{mfat2024}. Norwegian production of farmed Atlantic salmon increased from approximately 400,000 tonnes in 2000 to over 1.2 million tonnes in 2013, and has continued to grow, surpassing 1.5 million tonnes in recent years \citep{fiskeridir2024}. Norway currently accounts for roughly 50\% of global farmed Atlantic salmon production, making it the dominant supplier
in international markets \citep{ytrestoeyl2015}.

Salmon is Norway's second largest export commodity after oil and
gas, generating billions of Krones in annual export revenues and providing
employment in coastal communities that do not have many other economic
opportunities. The industry is geographically concentrated along Norway's
western and northern coastline, where farms are located in open-sea net
pen cages positioned in the cold, well-oxygenated waters of the fjords
\citep{fiskeridir2024}.

The sector is dominated by a small number of large companies, with some of the biggest being Mowi (formerly Marine Harvest), Lerøy Seafood, and SalMar. Each of these companies operate dozens of farm sites and also manage feed procurement, processing, and distribution systems. 

\section{Impact Mechanisms on Marine Biodiversity}

Despite how valuable this industry is for Norway's economy, the expansion of Norwegian salmon aquaculture has raised persistent concerns about its negative effects on the local marine biodiversity. The literature identifies three main impact mechanisms.

\paragraph{Feed and wild fish extraction:}
Farmed Atlantic salmon are carnivorous and require feed rich in protein and omega-3 fatty acids. These nutritional requirements have historically been met by using fishmeal and fish oil derived from wild-caught forage fish, small pelagic species such as anchovies, herrings, sand eels, and sprat, which occupy the lower levels of marine food webs. The issue is that extracting these species at a large enough scale to supply salmon feed for the whole industry reduces the availability of prey for their natural predators such as fish, seabirds, and marine mammals. This can also potentially have a cascading effect through the marine ecosystem structure \citep{tacon2008, oceanconservation}. The total mass of wild fish extracted to supply Norwegian salmon feed is the first biodiversity pressure we quantify in this report, denoted as \(W\).

Measuring the scale of this extraction is not straightforward. It depends on many factors: total salmon produced, how much feed is needed per kilogram of salmon produced, what share of the feed used consists of marine ingredients, and how much wild fish is needed to produce each kilogram of marine ingredient. The latter depending heavily on whether those ingredients come from dedicated forage fish catches or from fish-processing by-products. These four factors form the basis of our first IPAT decomposition, described in Section~4.

\paragraph{Sea lice and chemical treatments:}
Sea lice (otherwise known as \textit{Lepeophtheirus salmonis}) are naturally occurring ectoparasites of salmon. In the wild the density of host fish is typically low enough to keep infestations in check. However, when using sea pen aquaculture the high density of fish creates conditions that make it significantly easier for lice to spread \citep{nhm2024}. Beyond just the welfare implications for farmed fish, this matters for the nearby marine biodiversity because sea lice larvae are released into surrounding waters and can infest wild Atlantic salmon and sea trout migrating past the farms \citep{plos2018}. Wild salmon populations near Norwegian aquaculture sites are already under a lot of pressure, and additional lice-induced mortality can reduce both abundance and genetic diversity which is especially important for salmon \citep{nhm2024}.

To control sea lice, farms typically rely on a combination of chemical and non-chemical treatments. Chemical treatments can be divided into two categories, bath treatments and oral treatments. Bath treatments are where the fish are enclosed and exposed to an antiparasitic compound dissolved in seawater, and oral treatments are treatments in which the active substance is incorporated into the feed to be eaten by the salmon \citep{myhrjensen2020}. The environmental impacts of these two methods differ sharply, bath treatments release virtually all of the active substance directly into the fjord (approximately 100\% discharge), while the oral treatments are much lower due to the fact that 97\% of the active ingredients are metabolized before they reach the water (leaving an approximate 3\% discharge)\citep{myhrjensen2020}. Pyrethroids such as cypermethrin and deltamethrin, and the organophosphate azamethiphos are widely used in the bathing process and are toxic to non-target crustaceans, including European lobster (\textit{Homarus gammarus}), shrimp, and other invertebrates living near the farm sites \citep{samuelsen2020}. The mass of persistent antiparasitic substances discharged annually into Norwegian fjords is the second biodiversity pressure we quantify in this report, denoted as \(A_{out}\).

A third mechanism does exist but is not investigated in this report. The escape of farmed Atlantic salmon into the wild, is recognised in the scientific literature as a source of biodiversity pressure. More specifically it is when escaped salmon change the local genetic pool and compete with the wild salmon. Farmed salmon are typically bred for rapid growth and differ genetically from the wild populations. When they escape their net pens and breed with wild fish, the resulting offspring tend to have lower survival rates in natural environments. This reduces the fitness and long-term health/survivability of wild Atlantic salmon populations. This mechanism is not captured by either of our IPAT decompositions, as consistent national data on escape rates, river entry points, and number of escaped salmon are insufficient to construct a reliable third identity. It is nonetheless included here for completeness, and a decomposition framework for genetic impact is proposed as a direction for future work in Section 9.3.

There are definitely many other recognised pressures from salmon aquaculture but they all fall outside the scope of this analysis, including greenhouse gas emissions from boats, pollution from fish processing, and the accumulation of fecal matter from the salmon beneath farm sites. The last one is especially surprising in its scale, a 2026 report by the Sunstone Institute estimates that the discharge from Norwegian aquaculture is comparable in magnitude to the raw sewage of millions of people, representing yet another large pressure on fjord ecosystems \citep{sunstone2026, guardian2026}. We focus on wild fish extraction and persistent chemical discharge because these are the two pressures where consistent time-series data exists to construct a complete IPAT decomposition over our study period. We emphasize that this choice is  data-driven, not a judgment about the importance of each of these pressures.



\section{The Decoupling Framework}

Several previous studies have examined components of the relationship between
Norwegian salmon production and its environmental footprint.
\citep{ytrestoeyl2015} documented the decline in the marine
ingredient share of Norwegian salmon feed between 2000 and 2013, showing
that the share of fishmeal and fish oil in the feed fell from
approximately 65\% to 29\% over this time period. This appears to be largely driven by the substitution of marine proteins with plant-based alternatives \citet{tacon2008}. On the treatment side, \citet{myhrjensen2020} provides the
most comprehensive analysis to date of the trends in sea lice treatment
methods in Norwegian aquaculture, documenting the structural shift away
from bath treatments and towards oral and non-chemical alternatives
between 2012 and 2019.

Our contribution builds directly on this body of work. Rather than examining individual efficiency metrics in isolation, our work focuses on making them into two complete IPAT decompositions that allow us to investigate the net effect of all relevant factors to assess whether absolute decoupling has occurred, and to which specific drivers have had the biggest impact. Applying this framework to both pressures allows for a comparative analysis of the mechanisms and policies that have shaped each outcome.














%==================================================
% 3. IPAT DECOMPOSITION — THEORY AND METHODOLOGY
%==================================================
\chapter{Theory: IPAT equation}

Determining whether the growth of Norwegian salmon aquaculture has led to an increase in its environmental footprint requires creating an IPAT equation. In this section, we present the IPAT approach that underpins our analysis, explain how we have adapted it to the two biodiversity pressures we study, and describe the structure of each decomposition.

\section{IPAT 1 : Wild Fish Extraction}

Our first decomposition tracks $W$, which represents the total mass of wild fish extracted from the ocean annually to supply Norwegian salmon feed. $W$ is not directly reported in any official statistics but it is a derived quantity that we compute from the following identity:

\begin{equation}
    W = P \times \frac{F}{P} \times \frac{M}{F} \times \frac{W}{M}
\end{equation}

Each term captures a step in the chain from salmon production
to wild fish extraction:

\begin{table}[h]
\centering
\caption{IPAT 1 decomposition terms}
\begin{tabularx}{\textwidth}{c X c c}
\hline
\textbf{Symbol} & \textbf{Definition} & \textbf{Unit} & \textbf{Source} \\
\hline
$W$   & Total wild fish extracted for Norwegian salmon feed
      & t\,yr$^{-1}$ & Computed \\[4pt]
$P$   & Atlantic salmon harvested in Norway
      & t\,yr$^{-1}$ & Fiskeridir \\[4pt]
$F/P$ & Economic feed conversion ratio (eFCR)
      & kg\,kg$^{-1}$ & Nofima \\[4pt]
$M/F$ & Marine ingredient share of compound feed
      & kg\,kg$^{-1}$ & Nofima \\[4pt]
$W/M$ & Wild forage fish per kg of marine ingredient (FFDR\textsubscript{FM})
      & kg\,kg$^{-1}$ & Nofima \\
\hline
\end{tabularx}
\label{tab:ipat1}
\end{table}

The factor F/P measures how much feed is needed to produce one kilogram of salmon, known as the economic Feed Conversion Ratio or eFCR. A lower eFCR means less feed is required for each kilogram of fish, and therefore fewer marine resources are used per kilogram. The factor $M/F$ captures the marine ingredient share of that feed. This share can be reduced by substituting plant-based proteins and novel ingredients for fishmeal and fish oil.
The factor $W/M$ measures how much wild fish is needed to produce marine ingredients. It decreases when more of these ingredients come from fish by-products like trimmings and offcuts, instead of fish caught specifically for feeding salmon \citep{tacon2008, ytrestoeyl2015}.

\paragraph{Scope and period:}

The analysis covers Atlantic salmon only, and not rainbow trout, which
is also farmed in Norway. This is primarily due to data consistency:
the Nofima benchmark reports that provide our $F/P$, $M/F$, and $W/M$
values are based on feed data collected from the three major Norwegian
feed producers : BioMar, EWOS, and Skretting. 

The study period covers the years from 2000 to 2013. This window is determined by the availability of Nofima benchmark data. The reports provide point estimates for 2000, 2010, 2012, and 2013, which are the four years for which $F/P$, $M/F$, and $W/M$ can be derived with confidence. This window also captures the most dramatic phase of the transition in feed composition, during which the marine ingredient share fell from 64.6\% to 29.2\%.

\section{IPAT 2: Antiparasitic Chemical Discharge}

Our second decomposition tracks $A_{out}$, the annual mass of
persistent antiparasitic substances released into Norwegian fjords
as a result of sea lice treatments, used here as an indicator of
persistent chemical discharge pressure. It is decomposed into the five
 factors as follows:

\begin{equation}
    A_{out} = N_{salmon}
    \times \frac{N_{lice}}{N_{salmon}}
    \times \frac{Tr_{total}}{N_{lice}}
    \times \frac{Tr_{chem,pol}}{Tr_{total}}
    \times \frac{A_{out}}{Tr_{chem,pol}}
\end{equation}

where $\lambda \equiv N_{lice}/N_{salmon}$, $\tau \equiv Tr_{total}/N_{lice}$, 
$\phi \equiv Tr_{chem,pol}/Tr_{total}$, and $\delta \equiv A_{out}/Tr_{chem,pol}$ 
are defined in Table~\ref{tab:ipat2}.


\begin{table}[h]
\centering
\caption{IPAT 2 decomposition terms}
\begin{tabularx}{\textwidth}{c c X c}
\hline
\textbf{Factor} & \textbf{Symbol} & \textbf{Definition} & \textbf{Unit} \\
\hline
$N_{salmon}$
    & ---
    & Total number of Atlantic salmon harvested in Norway
    & fish\,yr$^{-1}$ \\[4pt]
$N_{lice}/N_{salmon}$
    & $\lambda$
    & Average number of adult female sea lice per fish,
    captures infestation intensity
    & lice\,fish$^{-1}$ \\[4pt]
$Tr_{total}/N_{lice}$
    & $\tau$
    & Treatment events per unit of lice burden, measures
    how aggressively the industry responds to a given
    infestation level regardless of treatment type
    & events\,lice$^{-1}$ \\[4pt]
$Tr_{chem,pol}/Tr_{total}$
    & $\phi$
    & Share of treatment events using persistent chemicals
    rather than mechanical, thermal, or biological
    alternatives, the key decoupling term
    & dimensionless \\[4pt]
$A_{out}/Tr_{chem,pol}$
    & $\delta$
    & Kilograms of persistent substance discharged per
    chemical treatment event. falls sharply as the
    treatment mix shifts from bath (100\% discharge)
    toward oral ($\sim$3\% discharge)
    & kg\,event$^{-1}$ \\
\hline
\end{tabularx}
\label{tab:ipat2}
\end{table}

The definition of $A_{out}$ used in our calculations
is:

\begin{equation}
    A_{out} = A_{\text{bath, excl.\,H}_2\text{O}_2} \times 1.0
                + A_{\text{oral}} \times 0.03
\end{equation}

Hydrogen peroxide (H$_2$O$_2$) is excluded because it degrades in
seawater with a half-life of less than one hour and does not persist
in the marine environment \citep{myhrjensen2020}. Including it would
cause it to represent over 99.9\% of total administered substance in
some years, completely masking the relevant signal from the actual
persistent compounds. Oral treatments are assigned a discharge factor
of 0.03 (3\%), reflecting the fact that approximately 97\% of orally
administered substances are metabolised by the fish before reaching
the fjord \citep{myhrjensen2020}.


\paragraph{Study period:}
This analysis covers 2012 to 2019, the period for which the treatment
data required to compute $Tr_{total}/N_{lice}$,
$Tr_{chem,pol}/Tr_{total}$, and $A_{out}/Tr_{chem,pol}$ are
available from the BarentsWatch Fish Health API \citep{barentswatch},
the Norwegian Directorate of Fisheries \citep{fiskeridir2024},
and the study by \citet{myhrjensen2020}. This period also captures
the most important transition in the Norwegian lice
treatment methods used by the industry. During this time we see the near-disappearance of persistent bath treatments and the rise of non-chemical alternatives.












%==================================================
% 4. DATA SOURCES AND ASSUMPTIONS
%==================================================
\chapter{Data Sources and Assumptions}

The two IPAT decompositions presented in this report rely on a small
number of carefully selected data sources, each corresponding to one
or more factors in the decomposition identities described in Section~4.
In this section, we describe each source, explain how it was used, and
state the key assumptions we make.

\section{Data Sources for IPAT 1}

\paragraph{Salmon production ($P$):}
Annual Atlantic salmon production data are taken from the official 
statistics of the Norwegian Directorate of Fisheries \citep{fiskeridir2024}. The dataset reports annual sales volumes broken down by county and species. The resulting annual series covers 2000 to 2013 and is expressed in tonnes of round weight.

\paragraph{Feed efficiency, marine ingredient share, and wild fish
dependency ($F/P$, $M/F$, $W/M$):}
The three efficiency factors are derived from the Nofima benchmark
reports on Norwegian salmon feed composition. Specifically, we use two
sources: \citet{ytrestoeyl2015} published in \textit{Aquaculture} (vol.
448, pp. 365--374), and Nofima Report 36/2014 by Ytrestøyl, Aas, and
Åsgård. These reports provide four point estimates during our timespan: 
2000, 2010, 2012, and 2013. The values used are summarised in
Table~\ref{tab:nofima}.

\begin{table}[h]
\centering
\caption{Nofima benchmark data for IPAT 1 efficiency factors}
\begin{tabular}{cccc}
\hline
\textbf{Year} & \textbf{$M/F$} & \textbf{$F/P$ (eFCR)} &
\textbf{$W/M$ (FFDR\textsubscript{FM})} \\
\hline
2000 & 0.646 & 1.50 & 1.80 \\
2010 & 0.414 & 1.20 & 1.10 \\
2012 & 0.307 & 1.15 & 0.70 \\
2013 & 0.292 & 1.15 & 0.70 \\
\hline
\end{tabular}
\label{tab:nofima}
\end{table}

The $W/M$ value is defined as the Forage Fish Dependency Ratio for
fishmeal (FFDR\textsubscript{FM}), which captures how much raw wild
catch is needed to produce one kilogram of usable marine ingredient,
excluding processing by-products.

\section{Data Sources for IPAT 2}

\paragraph{Production scale ($N_{\text{salmon}}$):}
$N_{\text{salmon}}$ is derived from $P$ (total Atlantic salmon 
produced in tonnes, same Fiskeridir source as IPAT 1) via:
\begin{equation}
    N_{\text{salmon}} = \frac{P \times 1000}{w}
\end{equation}
assuming a constant average slaughter weight of $w = 5.0$ kg/fish, 
which is a value consistent with those reported for Norwegian Atlantic salmon over this period \citep{ytrestoeyl2015}. $N_{\text{salmon}}$ is preferred over $P$ as the scale driver in 
IPAT 2 because the downstream factors (lice burden per fish 
($\lambda$), treatment events per lice ($\tau$), and chemical 
share ($\phi$))  are all defined at the individual fish level, 
making the number of fish the natural unit of analysis.



\paragraph{Lice counts, treatment events, and substance quantities:}
The data required to compute the four lice-related factors 
$N_{lice}/N_{salmon}$, $Tr_{total}/N_{lice}$,
$Tr_{chem,pol}/Tr_{total}$, and $A_{out}/Tr_{chem,pol}$ come from 
three sources. Lice counts and treatment event records are retrieved 
from the BarentsWatch Fish Health API \citep{barentswatch}, providing 
year-by-year national statistics on average adult female lice per fish 
and total treatment events by treatment type. Quantities of active 
substances administered by substance and delivery method are taken 
from \citet{myhrjensen2020}, which provides a consolidated and 
cross-validated dataset covering 1992 to 2017. Production data 
($N_{\text{salmon}}$) are derived from the same Fiskeridir source 
as IPAT~1 \citep{fiskeridir2024}.

\section{Key Assumptions}

\paragraph{Linear interpolation of Nofima benchmark data:}
The Nofima reports provide point estimates for only four years (2000,
2010, 2012, 2013) within our 2000 to 2013 window. To obtain annual
values for the intervening years, we apply piecewise linear interpolation
between consecutive points using the \texttt{approx()} function
in R. This is the most important assumption in IPAT~1.

Linear interpolation is justified here for two reasons. First, the
four benchmark points span the full study period and the transition
in feed composition was gradual and continuous, making linear
interpolation a reasonable approximation. Second, the endpoint
comparison, which is the basis of our main conclusion, relies
entirely on measured benchmark values for 2000 and 2013, and is
therefore not affected by the interpolation. The year-by-year
decomposition charts should be interpreted as indicative
rather than precise, particularly in the 2001 to 2009 period where
the only available benchmarks are 2000 and 2010.

\paragraph{Atlantic salmon only:}
As explained in Section~4, our IPAT~1 analysis covers Atlantic salmon
production only, excluding rainbow trout. This is necessary for internal
consistency with the Nofima benchmark data, which are based on Atlantic
salmon feed. Using total Norwegian salmonid production as $P$ while
applying salmon-specific efficiency factors would systematically
overestimate $F/P$.

\paragraph{National aggregation:}
In IPAT~2, we use national aggregate data on treatment events and
substance quantities, which cover all Norwegian salmon farms collectively.
We do not attempt to isolate the contribution of individual production
areas or farm sizes. This means that our results reflect national
averages and may mask regional variation in treatment intensity or
treatment mix. Farms in areas with higher lice pressure, or in zones
classified as red under the Traffic Light System, may show very different
patterns from the national aggregate.

\paragraph{Discharge factors for bath and oral treatments:}
We assign a discharge factor of 1.0 (100\%) to persistent bath treatments
and 0.03 (3\%) to oral treatments, following \citet{myhrjensen2020}.
These values represent the best available estimates in the scientific
literature but are themselves subject to uncertainty. The 3\% oral
discharge factor, in particular, is a population average that may vary
depending on the specific compound used, the feeding behaviour of the
fish, and water temperature. Sensitivity to this assumption is not
formally explored in our analysis, but we note that even if the oral
discharge factor were doubled to 6\%, the overall conclusion, that
$A_{out}$ declined by more than 90\% over the study period, would
remain unchanged, given the magnitude of the observed decline.

\paragraph{Exclusion of hydrogen peroxide:}
As stated in Section~4, H$_2$O$_2$ is excluded from our $A_{out}$
definition on the grounds that it degrades in seawater within less than
one hour and does not accumulate in the marine environment. This is an
ecologically motivated choice. Our indicator is designed to capture
persistent pollution, not total administered mass. Including H$_2$O$_2$
would change the absolute magnitude of $A_{out}$ dramatically but would
not change the direction or approximate magnitude of the trend, since
H$_2$O$_2$ usage and persistent chemical usage are driven by different
factors and follow different trajectories over our study period.









%==================================================
% 5. RESULTS — IPAT 1: WILD FISH EXTRACTION
%==================================================
\chapter{Results IPAT 1: Wild Fish Extraction}

This section presents the results of our first IPAT decomposition,
covering the period 2000 to 2013. We first document the growth in
salmon production, then trace the evolution of each efficiency factor,
and finally show how their combined effect produced a strong absolute
decoupling between production and wild fish extraction.

\section{Production Growth: The Scale Effect}

Between 2000 and 2013, Atlantic salmon production in Norway grew from
approximately 440,000 tonnes to 1,168,000 tonnes, an increase of 166\%.
Figure~\ref{fig:P} shows this trajectory.

\begin{figure}[H]
    \centering
    \includegraphics[width=0.85\textwidth]{Images/Atlantic_Salmon_production.png}
    \caption{Atlantic salmon production in Norway, 2000--2013 (tonnes
    round weight). Norwegian Directorate of Fisheries \citep{fiskeridir2024}}
    \label{fig:P}
\end{figure}

This production growth is the scale effect in our decomposition. All
else being equal, a 166\% increase in $P$ would produce a 166\% increase in
$W$. The fact that $W$ instead declined sharply over this period as we show below means that the three efficiency factors must have more than fully compensated for this scale effect.

\section{Evolution of the Three Efficiency Factors}

Figure~\ref{fig:interp} shows the interpolated annual time series for
the three efficiency factors $F/P$, $M/F$, and $W/M$, with the four
Nofima benchmark points shown as open circles.

\begin{figure}[H]
    \centering
    \includegraphics[width=\textwidth]{Images/Ipat1_Interpolation_factors.png}
    \caption{Piecewise linear interpolation of the three IPAT~1
    efficiency factors between Nofima benchmark points (2000, 2010,
    2012, 2013). Open circles indicate measured benchmark values.
    Sources: \citet{ytrestoeyl2015} and \citet{nofima2014}.}
    \label{fig:interp}
\end{figure}

\paragraph{Feed conversion ratio ($F/P$):}
The economic feed conversion ratio fell from 1.50 in 2000 to 1.15 in
2013, a reduction of 23\%. This improvement reflects advances in feed
formulation, feeding management, and fish genetics that have allowed
farms to produce more salmon per kilogram of feed consumed. In practical
terms, a lower $F/P$ means that less feed and therefore fewer marine
ingredients is needed to produce each tonne of salmon. The decline
in $F/P$ was relatively steady over the period, with most of the
improvement occurring in the 2000 to 2010 window, and the ratio
stabilising at 1.15 from 2012 onwards.

\paragraph{Marine ingredient share ($M/F$):}
The marine ingredient share of Norwegian compound salmon feed fell from
0.646 (64.6\%) in 2000 to 0.292 (29.2\%) in 2013, a reduction of 55\%.
This reflects a change in the composition of salmon feed
over this period: plant-based proteins, primarily soy protein
concentrate and rapeseed meal, progressively replaced fishmeal,
while vegetable oils partially replaced fish oil. The sharpest decline
occurred between 2000 and 2010, after which the rate of substitution
slowed as the marine ingredient share approached a practical minimum
consistent with the nutritional requirements of Atlantic salmon.

\paragraph{Wild fish dependency ratio ($W/M$):}
The ratio of wild forage fish required per kilogram of marine ingredient
fell from 1.80 in 2000 to 0.70 in 2013, a reduction of 61\%.
This decline reflects a growing share of marine ingredients sourced
from fish-processing by-products, including trimmings and offcuts from fish
processing plants that would otherwise be discarded rather than
from dedicated catches of forage fish such as anchovies or sand eels.
Since by-products do not require a dedicated wild fish catch, their
use reduces $W/M$ without requiring any reduction in $M/F$. The
sharpest decline occurred between 2010 and 2012, when the ratio fell
from 1.10 to 0.70, suggesting that this shift accelerated markedly
in the early 2010s.

\section{Wild Fish Extraction: Absolute Decoupling}

Figure~\ref{fig:Decoupling_1} shows the annual trajectories of $W$ (total wild
fish extracted) and $P$ (salmon produced), both expressed in million
tonnes, over the full 2000 to 2013 period.

\begin{figure}[H]
    \centering
    \includegraphics[width=1\textwidth]{Images/Absolut_decoupling.png}
    \caption{Wild fish extraction $W$ versus salmon production $P$,
    2000--2013 (million tonnes).}
    \label{fig:Decoupling_1}
\end{figure}

While $P$ grew continuously from 0.44 Mt to 1.17 Mt, $W$ moved in the opposite direction, falling from 0.77 Mt in 2000 to 0.27 Mt in 2013. This represents a decline of 64\% in absolute wild fish extraction over a period in which the industry more than doubled its output. This is a clear case of \textbf{absolute decoupling}, the environmental pressure $W$ decreased in absolute terms even as the scale of the industry expanded substantially.

\section{Ln-Additive Growth Factor Decomposition}

To quantify the contribution of each factor to the total change in $W$,
we apply a ln-additive decomposition. Figure~\ref{fig:waterfall1} shows the results as a waterfall chart, where each bar represents the ln-contribution of one factor to the
total ln-change in $W$ between 2000 and 2013.

\begin{figure}[H]
    \centering
    \includegraphics[width=0.85\textwidth]{Images/Waterfall_chart_factors_IPAT_1.png}
    \caption{Waterfall decomposition of $\Delta \ln W$, 2000--2013.}
    \label{fig:waterfall1}
\end{figure}

The decomposition shows the following ln-contributions:

\begin{itemize}
    \item \textbf{$P$ (production scale)}: $+0.98$ ln-points,
    corresponding to the 166\% increase in salmon output. This is the
    only factor pushing $W$ upward.
    \item \textbf{$F/P$ (feed conversion ratio)}: $-0.27$ ln-points,
    corresponding to the 23\% improvement in feed efficiency. This is
    the smallest of the three mitigating factors.
    \item \textbf{$M/F$ (marine ingredient share)}: $-0.79$ ln-points,
    corresponding to the 55\% reduction in the marine fraction of feed.
    This is the second largest mitigating factor.
    \item \textbf{$W/M$ (wild fish dependency)}: $-0.94$ ln-points,
    corresponding to the 61\% reduction in wild fish needed per kilogram
    of marine ingredient. This is the largest single mitigating factor.
    \item \textbf{$W$ (net result)}: $-1.03$ ln-points, corresponding
    to the observed 64\% decline in total wild fish extraction.
\end{itemize}

The three efficiency gains sum to $-2.00$ ln-points, more than twice
the scale effect of $+0.98$. Production growth alone would have more
than doubled wild fish extraction but the combined improvements in feed
efficiency and ingredient sourcing reduced it by nearly two thirds.

\section{Summary of IPAT 1 Results}

Table~\ref{tab:ipat1results} summarises the key figures for the full
2000--2013 period.

\begin{table}[H]
\centering
\caption{Ln-additive growth factor decomposition, IPAT 1 (2000--2013)}
\begin{tabular}{lrrrr}
\hline
\textbf{Factor} & \textbf{2000} & \textbf{2013} &
\textbf{\% change} & \textbf{ln contribution} \\
\hline
$P$ (production, t)       & 439,874  & 1,168,324 & $+166\%$ & $+0.977$ \\
$F/P$ (eFCR)               & 1.50     & 1.15      & $-23\%$  & $-0.266$ \\
$M/F$ (marine share)      & 0.646    & 0.292     & $-55\%$  & $-0.794$ \\
$W/M$ (FFDR\textsubscript{FM}) & 1.80 & 0.70      & $-61\%$  & $-0.944$ \\
\hline
$W$ (wild fish, t)        & 767,228  & 274,626   & $-64\%$  & $-1.027$ \\
\hline
\end{tabular}
\label{tab:ipat1results}
\end{table}

The results confirm absolute decoupling between salmon production and wild fish extraction over the 2000 to 2013 period. The 64\% decline in $W$ despite a 166\% increase in $P$ is driven by three simultaneous and mutually reinforcing improvements: a 23\% gain in feed efficiency, a 55\% reduction in the marine ingredient share of feed, and a 61\% reduction in the wild fish dependency of marine ingredients. Together, these three factors reduced the wild fish extraction intensity of Norwegian salmon production by approximately 86\% over the study period, more than enough to counteract the effect of doubled output.
















%==================================================
% 6. RESULTS — IPAT 2: ANTIPARASITIC CHEMICAL DISCHARGE
%==================================================
\chapter{Results IPAT 2: Antiparasitic Chemical Discharge}

This section presents the results of our second IPAT decomposition,
covering the period 2012 to 2019. We first document the overall
decoupling between salmon production and chemical discharge, then
trace the evolution of each of the five factors, and finally quantify
their individual contributions through a ln-additive waterfall
decomposition.

\section{Production Scale and the Baseline}

Over the 2012 to 2019 study period, Atlantic salmon production in Norway
remained broadly stable, growing modestly from approximately 1,232,000
tonnes in 2012 to 1,364,042 tonnes in 2019, an increase of 10.7\%.
This is in sharp contrast to the dramatic production growth documented
in IPAT~1. The 2012 to 2019 period falls within the post-peak phase of
the industry's rapid expansion, during which the Traffic Light System
and other regulatory constraints began to slow the rate of capacity
growth. The relatively flat production trajectory over this period is
important context for interpreting the IPAT~2 results because unlike in
IPAT~1, the scale effect here is small, meaning that the observed
changes in $A_{out}$ are driven almost entirely by changes in the
lice management factors rather than by changes in production volume.

\section{Absolute Decoupling: $A_{out}$ vs Production}

Figure~\ref{fig:decoupling2} shows the trajectories of $A_{out}$ and
salmon production $P$, both indexed to their 2012 values (2012 = 100),
over the full 2012 to 2019 period.

\begin{figure}[h]
    \centering
    \includegraphics[width=0.85\textwidth]{Images/Absolut_Decoupling_IPAT_2.png}
    \caption{Absolute decoupling between persistent chemical discharge
    $A_{out}$ and salmon production $P$.
    Source: \citet{myhrjensen2020}, \citet{barentswatch}, and \citet{fiskeridir2024}.}
    \label{fig:decoupling2}
\end{figure}

The decline is large. While production grew modestly, $A_{out}$ collapsed from 4,484 kg in 2012 to 212 kg in 2019, a decline of 95.3\%. The two series move in opposite directions throughout the period, with the divergence accelerating sharply after 2015. This is a case of absolute decoupling which is evident because the environmental pressure $A_{out}$ fell to less than 5\% of its 2012 level even as the industry continued to grow.

Figure~\ref{fig:Aout_abs} shows $A_{out}$ in absolute terms (kg
per year), which makes the magnitude of the decline more concrete.

\begin{figure}[H]
    \centering
    \includegraphics[width=0.85\textwidth]{Images/A_out_absolut_term.png}
    \caption{Persistent antiparasitic substances discharged into
    Norwegian fjords ($A_{out}$), kg per year, 2012--2019.
    Source: \citet{myhrjensen2020}, \citet{barentswatch}, 
and \citet{fiskeridir2024}.}
    \label{fig:Aout_abs}
\end{figure}

Several features of the absolute $A_{out}$ series are worth noting.
First, discharge was not monotonically declining, it rose from 4,484
kg in 2012 to a peak of 5,186 kg in 2014 before beginning its sharp
decline. This temporary increase reflects a period of high lice
pressure in 2014, which drove up the number of chemical treatment
events before the structural shift towards non-chemical alternatives
had fully taken hold. Second, the sharpest single-year drop occurs
between 2015 and 2016, when $A_{out}$ fell from 4,364 kg to 1,638
kg, a reduction of 62\% in a single year marking the point at which the transition away from persistent bath treatments became decisive at the national scale. By 2017, $A_{out}$ had fallen to 293 kg, and it remained below 220 kg in both 2018 and 2019.

\section{IPAT 2 Factor Time Series}

Figure~\ref{fig:factors2} shows the individual time series for each
of the five factors in the IPAT~2 decomposition, allowing us to
trace the trajectory of each driver independently over the study
period.

\begin{figure}[h]
    \centering
    \includegraphics[width=\textwidth]{Images/IPAT_2_Factors.png}
    \caption{Individual factor time series for the IPAT~2
    decomposition, 2012--2019. The identity is $A_{out} = N_{\text{salmon}}
    \times \lambda \times \tau \times \phi \times \delta$.
    Source: \citet{myhrjensen2020}, \citet{barentswatch}, 
and \citet{fiskeridir2024}.}
    \label{fig:factors2}
\end{figure}

\paragraph{Lice pressure ($N_{lice}/N_{salmon}$):}
The mean number of adult female lice per fish fluctuated considerably
over the study period, reflecting the natural variability of sea lice
infestation levels. It was highest at the start of the study period
in 2012 (0.214 lice per fish) and rose again to a secondary peak
in 2014 (0.207 lice per fish), which corresponds to the temporary
peak in $A_{out}$ noted above. It was lowest in 2018 (0.142 lice
per fish). Over the full period, this
factor declined modestly, contributing $-0.30$ ln-points to the
total change in $A_{out}$. Improvements in lice pressure
management were therefore present but were not the primary driver
of the decline in discharge.

\paragraph{Treatment intensity ($Tr_{total}/N_{lice}$):}
The number of treatment events deployed per unit of lice burden
fluctuated over the period, with a notable increase between 2015
and 2017 before declining again in 2018 to 2019. Its net contribution
over the full period is $-0.22$ ln-points. This relatively small
figure reflects the fact that, while farms were responding more
intensively to a given lice burden in some years, the nature of
those treatments was shifting simultaneously towards less polluting
methods, a change captured by the two following factors rather
than by this one.

\paragraph{Chemical treatment share ($Tr_{chem,pol}/Tr_{total}$):}
The share of total treatment events involving persistent chemical
compounds fell from approximately 50\% in 2012 to 27\% in 2019,
contributing $-0.60$ ln-points to the total change in $A_{out}$.
This reflects the growing use of non-chemical delousing methods,
mechanical treatments such as hydrolicer and snorkel systems, thermal
treatments using warm water, and biological control using cleaner
fish (lumpfish and wrasse), which progressively displaced chemical
treatments as the industry's primary response to sea lice. The
decline was particularly sharp between 2016 and 2017, when the
number of non-chemical treatment events more than doubled as
mechanical and thermal technologies became widely available at
commercial scale.

\paragraph{Discharge per chemical treatment event
($A_{out}/Tr_{chem,pol}$):}
This factor fell from 1.72 kg per event in 2012 to 0.22 kg per
event in 2019, contributing $-2.04$ ln-points to the total change
in $A_{out}$. This is by far the largest single contribution
in the decomposition, accounting for approximately 67\% of the
total observed decline. Its decline is explained by a shift within
the chemical treatment category itself: bath treatments, which
release approximately 100\% of the active substance directly into
the surrounding water, became progressively rarer, while oral
treatments, which result in only approximately 3\% of the
administered substance reaching the fjord, became more common.
The result is that even when a farm resorts to a chemical
treatment, the amount of persistent substance discharged per
event is now a fraction of what it was in 2012.

\section{Waterfall Decomposition}

Figure~\ref{fig:waterfall2} shows the ln-additive waterfall
decomposition of the total change in $A_{out}$ between 2012 and
2019.

\begin{figure}[H]
    \centering
    \includegraphics[width=0.85\textwidth]{Images/IPAT_2_Waterfall.png}
    \caption{Waterfall decomposition of $\Delta \ln A_{out}$,
    2012--2019. 
    Source: \citet{myhrjensen2020}, \citet{barentswatch}, 
and \citet{fiskeridir2024}.}
    \label{fig:waterfall2}
\end{figure}

The five ln-contributions are as follows:

\begin{itemize}
    \item \textbf{$N_{salmon}$ (production scale)}: $+0.10$,
    the only factor pushing $A_{out}$ upward, reflecting the
    modest 10.7\% production growth over the period.

    \item \textbf{$N_{lice}/N_{salmon}$ (lice pressure)}: $-0.30$,
    reflecting a moderate improvement in average lice levels per
    fish over the study period.

    \item \textbf{$Tr_{total}/N_{lice}$ (treatment intensity)}:
    $-0.22$, reflecting a modest net reduction in treatment events
    per unit of lice burden.

    \item \textbf{$Tr_{chem,pol}/Tr_{total}$ (chemical treatment
    share)}: $-0.60$, reflecting the significant shift towards
    non-chemical delousing methods after 2015.

    \item \textbf{$A_{out}/Tr_{chem,pol}$ (discharge per
    chemical event)}: $-2.04$, the dominant driver, reflecting
    the structural shift from bath to oral chemical treatments
    within the chemical treatment category.

    \item \textbf{$A_{out}$ (net result)}: $-3.05$,
    corresponding to the observed 95.3\% decline in total
    persistent discharge over 2012--2019.
\end{itemize}

\section{Treatment Mix Structural Shift}

Figure~\ref{fig:treatmix} shows the full treatment mix across all
intervention types for each year of the study period, expressed as
total events per year. This chart illustrates the underlying structural
changes in industry practice that drove the two dominant factors in
our decomposition $Tr_{chem,pol}/Tr_{total}$ and
$A_{out}/Tr_{chem,pol}$ by showing directly how the
composition of treatment events evolved over time.

\begin{figure}[H]
    \centering
    \includegraphics[width=0.85\textwidth]{Images/Treatment_Mix_shift_IPAT_2.png}
    \caption{Treatment mix across all intervention types (events per
    year), 2012--2019.
    Source: \citet{myhrjensen2020}, \citet{barentswatch}, 
and \citet{fiskeridir2024}.}
    \label{fig:treatmix}
\end{figure}

Several features of this chart deserve attention. The most important
is the near-complete disappearance of persistent bath treatments,
azamethiphos and pyrethroid compounds applied by enclosing the net
pen after 2016. From over 2,100 bath treatment events in 2012,
this category fell to fewer than 200 events by 2018. This structural
change is the primary driver of the decline in
$A_{out}/Tr_{chem,pol}$: as bath treatments were replaced by
oral and non-chemical alternatives, the average discharge per
chemical event collapsed from 1.72 kg in 2012 to 0.22 kg in 2019.

The rise of non-chemical treatments is equally visible. mechanical
and thermal delousing events, which were rare in 2012, grew to
represent the majority of all treatment events by 2017--2018,
driving the decline in $Tr_{chem,pol}/Tr_{total}$ documented above.
The cleaner fish category shows a large apparent decline in
2018 to 2019, falling from over 4,000 events in 2017 to near zero
by 2019. This drop does not reflect an actual cessation of cleaner
fish use, but rather a reclassification in the reporting system
following a welfare crisis in 2017--2018 that prompted regulatory
scrutiny of lumpfish and wrasse mortality rates
\citep{myhrjensen2020}.

\section{Summary of IPAT 2 Results}

The results of our second IPAT decomposition confirm absolute decoupling between salmon production and persistent antiparasitic discharge in Norway over the 2012 to 2019 period. $A_{out}$ fell by 95.3\% while production grew by 10.7\%. These figures are mass-based: they measure kilograms of persistent substance discharged, not toxicity-weighted ecological impact. A 95.3\% reduction in discharged mass does not imply a proportional reduction in ecological harm, since the substances involved differ substantially in their toxicity to non-target marine organisms (see Section 8.2 for full discussion). The dominant driver is $A_{out}/Tr_{chem,pol}$ ($-2.04$ ln-points), capturing the structural shift from bath to oral chemical treatments. A secondary but significant contribution comes from $Tr_{chem,pol}/Tr_{total}$ ($-0.60$ ln-points), reflecting the broader shift in the treatment mix away from chemical interventions altogether. The factors
$N_{lice}/N_{salmon}$ and $Tr_{total}/N_{lice}$ made smaller but consistent contributions of $-0.30$ and $-0.22$ ln-points respectively, while production growth exerted only a minor upward pressure of $+0.10$ ln-points throughout the period.
















%==================================================
% 3. REGULATORY FRAMEWORK
%==================================================
\chapter{Regulatory Framework}

The results from Sections~5 and~6 raise a natural question: why did the efficiency factors evolve the way they did? Why did $M/F$ and $W/M$ fall consistently between 2000 and 2013? Why did $\delta$ collapse after 2015, and why did $\phi$ decline sharply after 2016? To answer these questions, we turn to the Norwegian regulatory framework governing the salmon aquaculture industry.


\section{Licensing System and Maximum Allowed Biomass}


Access to Norwegian salmon aquaculture is controlled through a site licensing system administered by the Norwegian Directorate of Fisheries (\textit{Fiskeridirektoratet}). Each licence specifies a \textit{Maximum Allowed Biomass} (MAB), i.e. the maximum tonnage of live fish that may be held at the site at any point in time.

Total production capacity expanded over time through successive licence allocations, which explains how salmon output grew from approximately 400,000 tonnes in 2000 to over 1.2 million tonnes in 2013, the 166\% increase that constitutes the scale effect in IPAT~1. Crucially, the MAB system did not constrain production growth during this period. The absolute decoupling we document wild fish extraction $W$ fell by 64\% despite this production increase, is therefore not the result of supply-side caps, but entirely attributable to improvements in $F/P$, $M/F$, and $W/M$, driven by feed industry innovation and market incentives.

\section{Sea Lice Thresholds and Monitoring Obligations}

Norwegian law sets mandatory limits on sea lice infestation at farm sites. Farms must not exceed a mean of 0.5 adult female lice per fish during most of the year \citep{nes2025}, and 0.2 during the spring smolt migration period \citep{nes2025, fni2023}. These limits are legally binding: a farm that exceeds them must immediately initiate treatment.

Farms are also required to count lice weekly on a minimum of 20 fish per cage, covering at least half of the cages on site \citep{plos2018}. These counts are reported to the Norwegian Food Safety Authority (\textit{Mattilsynet}) and compiled into national databases accessible through the BarentsWatch portal \citep{barentswatch}.

These rules connect directly to our IPAT~2 decomposition. The lice threshold caps $\lambda$, the mean number of adult female lice per fish, ensuring it cannot rise far above the legal limit without triggering a response. The mandatory treatment requirement shapes $\tau$, the number of treatment events per unit of lice burden, which is why $\tau$ remained stable over our study period despite the overall decline in chemical discharge. %Importantly, the regulation requires that treatment occurs, but does not specify which type, leaving open the choice between chemical bath treatments, oral treatments, and non-chemical alternatives. This is precisely the dimension captured by $\phi$ and $\delta$.

\section{The Traffic Light System}

In 2017, the Norwegian Ministry of Trade, Industry and Fisheries introduced the Traffic Light System (\textit{trafikklyssystemet}, TLS) \citep{nes2025, seafoodsource2024}. The system divides the Norwegian coastline into 13 production areas, each assessed annually by the Institute of Marine Research (IMR). Based on the estimated mortality of wild salmon smolts attributable to sea lice, each area is assigned a colour. Green areas may expand licensed capacity by 6\%, yellow areas are frozen, and red areas must reduce capacity by 6\%.

This creates a direct financial incentive to keep lice levels low, since a red classification triggers a lasting reduction in permitted biomass. At the same time, farms were already required to treat whenever lice counts exceeded the legal threshold, regardless of the method chosen. In the context of IPAT~2, the TLS reinforced and locked in the ongoing decline in $\phi$, the share of chemical treatments, which had already begun before the system formally came into force. As pyrethroid resistance grew \citep{sciencedirect2020b} and the availability of non-chemical alternatives expanded, farms shifted toward mechanical, thermal, and biological delousing methods particularly after 2015--2016; the TLS then strengthened this trajectory by creating direct financial incentives to sustain it.

\section{Regulation of Antiparasitic Treatments}

All antiparasitic treatments must be prescribed by a veterinarian and are regulated under Norwegian pharmaceutical law. The authorised persistent chemical substances include bath-applied compounds, azamethiphos, cypermethrin and deltamethrin and orally administered compounds such as emamectin benzoate, diflubenzuron and teflubenzuron \citep{samuelsen2020}.

The environmental impact of these two delivery modes differs fundamentally. Bath treatments discharge virtually all of the active substance directly into the surrounding fjord (discharge fraction $\approx$ 100\%), whereas oral treatments are mostly metabolised by the fish before reaching the water (discharge fraction $\approx$ 3\%) \citep{myhrjensen2020}. This is the mechanistic foundation of $\delta$, as the treatment mix shifts from bath to oral delivery, $\delta$ falls sharply even if the total number of chemical treatments stays constant.

The shift away from bath treatments reflects several forces acting simultaneously. Pyrethroid resistance documented from the mid-2000s \citep{sciencedirect2020b}, made cypermethrin and deltamethrin less effective, pushing farms toward alternatives. This coincided with the rapid expansion of biological lice control using cleaner fish, primarily lumpfish (\textit{Cyclopterus lumpus}) and ballan wrasse (\textit{Labrus bergylta}), with approximately 50 million deployed annually in Norway by 2019 \citep{barrett2020}. The near-disappearance of persistent bath treatments after 2016 therefore reflects regulatory thresholds, TLS incentives, pyrethroid resistance, and the availability of non-chemical alternatives acting together.

\section{Feed Sustainability Policy}


Unlike sea lice regulation, the reduction in marine ingredient content in Norwegian salmon feed was not driven by direct state regulation. No Norwegian law has mandated specific limits on the fishmeal or fish oil content of aquaculture feed. Instead, the decline in $M/F$ and $W/M$ resulted from market incentives, voluntary industry commitments, and third-party certification standards.

The most significant of these is the Aquaculture Stewardship Council (ASC) standard, which requires that marine ingredients in certified feed come from sustainably managed fisheries. The three dominant feed producers, BioMar, Skretting, and Cargill/EWOS, have voluntarily reduced their dependence on wild-sourced marine ingredients, investing in alternative protein sources such as soy, rapeseed, insect meal, and single-cell proteins \citep{frontiers2024}.

A parallel driver of the decline in $W/M$ has been the growing use of fish-processing by-products as a source of marine ingredients. Since by-products do not require dedicated wild fish catches, their use reduces the wild fish dependency ratio without necessarily reducing the marine ingredient share \citep{ytrestoeyl2015}.

The absence of direct feed regulation means that the decoupling in IPAT~1 is primarily a market and technology story in contrast with IPAT~2.
















%==================================================
% 8. DISCUSSION
%==================================================
\chapter{Discussion}

%The two IPAT decompositions presented in Sections~5 and~6, and the
%regulatory analysis of Section~7, together provide a detailed picture
%of how and why the environmental footprint of Norwegian salmon
%aquaculture evolved over our two study periods. In this section, we
%step back from the individual results and examine what they mean
%together. We compare the two decouplings, evaluate the limits of our
%analytical approach, and situate our findings within the broader
%literature on decoupling and environmental governance.

\section{Comparing the Two Decouplings}


In IPAT~1, the decoupling is driven by feed industry
innovation. As documented in Section~7.5, the decline in $M/F$ and
$W/M$ was not mandated by direct regulation of feed composition. No
Norwegian law required feed producers to reduce their fishmeal content
or to source a specific share of marine ingredients from by-products.
The drivers were instead a combination of economic incentives,
fishmeal and fish oil are expensive commodities, and reducing their
share lowers feed costs, market pressure from certification
standards such as the ASC, and sustained investment in feed R\&D by
the major producers. The role of the state was indirect, it created
a monitoring and transparency framework that made sourcing practices
visible, and it supported the certification infrastructure that gave
market actors an incentive to improve. But the actual improvement in
feed composition was primarily a market and technology story.

In IPAT~2, the decoupling required a more active regulatory hand.
As documented in Sections~7.2, 7.3, and 7.4, three forces converged
to produce the structural shift in treatment practice that our
decomposition identifies as the primary driver of the 95\% decline
in $A_{out}$. First, the pre-existing mandatory sea lice thresholds
created the obligation to treat, and pyrethroid resistance combined
with the growing availability of mechanical, thermal, and biological
alternatives drove the initial shift away from persistent bath
treatments from 2015--2016 onward. Second, the Traffic Light System,
introduced in 2017, reinforced this trajectory by tying production
capacity directly to lice outcomes, creating stronger financial
incentives to maintain and deepen the transition to non-chemical
methods. In this case, the state played a much more direct role than
in IPAT~1, and the regulatory pressure of the pre-existing lice
thresholds and, subsequently, the TLS were essential in sustaining
the near-complete elimination of persistent chemical discharge at the
pace observed.

\section{Limitations of the IPAT Models}

\paragraph{Single-impact focus:}
Each of our IPAT identities captures one environmental pressure in
isolation. A reduction in $W$ does not guarantee an overall improvement
in the environmental footprint of salmon feed, the substitution of
fishmeal with soy protein concentrate raises questions about
deforestation and land use in soy-producing regions, and the substitution
with algae or insect meal involves its own resource and energy costs.
Similarly, a reduction in $A_{out}$ does not capture the environmental
costs of the non-chemical treatments that replaced chemical ones:
mechanical and thermal delousing can cause significant fish stress and
mortality, and the large-scale production of cleaner fish has its own
feed and welfare implications. As \citet{coolchoices2024} note in the
context of refrigerant regulation, a decrease in the studied impact is
not always synonymous with an overall improvement, as there may be an
increase in another impact not captured in the equation. Our
decompositions are subject to exactly this limitation.

\paragraph{Interpolation uncertainty in IPAT~1:}
The efficiency factors $F/P$, $M/F$, and $W/M$ are available only for
four benchmark years within the 2000 to 2013 window. The annual values
used in our decomposition are obtained by linear interpolation between
these points. This is a reasonable approximation for a period of
gradual and continuous change, but it means that the year-by-year
decomposition charts in Section~5 should be interpreted as indicative
rather than precise. The endpoint comparison (2000 vs 2013) is based
entirely on measured benchmark values and is therefore robust to this
assumption, but intermediate values carry interpolation uncertainty,
particularly in the 2001 to 2009 period where the only available benchmark
points are 2000 and 2010.

\paragraph{Mass-based rather than toxicity-weighted chemical impacts:}

A further limitation concerns the definition of $A_{out}$ in IPAT~2.
Our decomposition measures the mass of persistent antiparasitic
substances discharged into the marine environment, rather than a
toxicity-weighted ecological impact indicator. This means that the
model implicitly treats one kilogram of azamethiphos, deltamethrin,
cypermethrin, emamectin benzoate, diflubenzuron, and teflubenzuron
as equivalent, even though these substances differ substantially in
toxicity, persistence, and ecological effects on non-target marine
organisms. A more complete framework would weight each substance using
ecotoxicity characterization factors from life-cycle assessment (LCA)
studies or ecotoxicological literature.

We do not implement such toxicity weighting here because consistent
substance-specific ecotoxicity factors and environmental fate data
were not available across all compounds and years in our study period.
As a result, $A_{out}$ should be interpreted as a measure of persistent
chemical discharge pressure rather than as a direct measure of total
ecological damage. This significantly limits the strength of the
ecological conclusions that can be drawn from IPAT~2.

\paragraph{National aggregation:}
Both decompositions use national aggregate data, which may mask
substantial regional variation. Sea lice pressure, treatment intensity,
and treatment mix vary considerably across Norway's 13 production areas,
reflecting differences in oceanographic conditions, farm density, and
regulatory classification under the TLS. A farm in a red-classified
area may show very different factor trajectories from one in a green
area. Similarly, the marine ingredient share of feed may vary across
feed producers and farm operators. Our national averages are meaningful
for assessing aggregate trends, but they should not be used to draw
conclusions about individual farms or regions.


\paragraph{Scope of the biodiversity impact:}

A fourth caveat concerns scope. The two pressures we quantify, wild fish extraction and chemical discharge, are those for which the data required to construct a full decomposition identity were available. Escaped farmed salmon represent a third recognised mechanism, operating through genetic introgression and competitive displacement of wild conspecifics, but consistent national data on escape rates and subsequent river entry rates over our study periods are insufficient to support a comparable quantitative analysis. The absence of a third decomposition does not imply that this pressure is negligible: it reflects a data limitation rather than an ecological judgment. Whether absolute decoupling has occurred on this dimension remains an open empirical question, and the data requirements for answering it are outlined in Section~9.3. More broadly, recent research has further documented the scale of eutrophication pressure from organic waste deposition beneath farm sites. It is a pressure that our framework does not capture \citep{sunstone2026, guardian2026} underscoring that the two pressures we quantify represent only a subset of the industry's total environmental footprint.















%==================================================
% 9. RECOMMENDATIONS AND NEXT STEPS
%==================================================
\chapter{Recommendations and Next Steps}

\section{Maintaining Decoupling on Wild Fish Extraction}

The 64\% decline in wild fish extraction documented in IPAT~1 is an impressive achievement, but it should not be taken as a guarantee of continued improvement.

First, Norwegian authorities should consider establishing mandatory targets for the share of marine ingredients sourced from by-products or certified sustainable fisheries, rather than relying solely on voluntary commitments and third-party certification. As production continues to grow, by-product supply chains are finite and may not keep pace with demand, which would push $W/M$ back up.

Second, clear policy targets should be set for the marine ingredient share $M/F$ for example, a cap of 15\% marine ingredients in compound feed by a given year to give feed producers a long-term signal and incentivise continued investment in alternative protein sources such as insect meal, single-cell proteins, and algae-derived omega-3 oils.

Third, public research funding bodies and the major feed producers should jointly assess the remaining margin for feed conversion ratio improvement, since even modest gains at the current scale of production translate into substantial reductions in absolute feed demand.

\section{Consolidating Decoupling on Chemical Discharge}

The 95.3\% decline in persistent antiparasitic discharge documented in IPAT~2 is fragile. The near-disappearance of persistent bath treatments could be reversed if lice pressure increases significantly or if resistance to current treatment alternatives develops.

To address this, we recommend two targeted interventions. First, the Traffic Light System should be reinforced with a direct financial charge on sea lice emissions, a lice tax that would make the environmental cost of infestations visible to individual operators rather than shared across entire production zones. Such an instrument would directly target $\lambda$ and $\phi$, the two factors our decomposition identifies as having the most room for further improvement.

Second, investment in preventive lice management should be prioritised through physical barriers such as lice skirts and snorkel systems, and continued investment in sea lice vaccines \citep{fishsite2024}. The current approach has largely solved the discharge problem by changing how lice are treated, but has not yet solved the infestation problem itself.

\section{Directions for Future Research}

Several extensions would increase the value and policy relevance of this work.

Extending IPAT~1 to more recent years would allow us to assess whether the absolute decoupling in wild fish extraction has been maintained as production has grown beyond 2013, and whether the rate of improvement in $M/F$ and $W/M$ has slowed as the marine ingredient share approaches its practical minimum.

Extending IPAT~2 beyond 2019 would allow us to assess whether the dramatic decline in $A_{out}$ has been maintained, and whether any new treatment categories have emerged that require updating the decomposition framework.

Constructing a third IPAT decomposition for genetic impact from fish escapes, expressing total introgression as a function of production scale, escape rate per unit of biomass, proportion of escapees entering rivers, and introgression rate per spawning escapee, would allow a more comprehensive assessment of whether absolute decoupling has been achieved across all main biodiversity pressures.

Finally, a disaggregated analysis at the level of the 13 TLS production areas would allow us to test whether the national patterns we document are consistent across regions, and to examine whether the TLS produced different outcomes in red, yellow, and green zones.

A further strategic question concerns the long-term sustainable level
for both pressures studied, and whether it can realistically be reached
using the methods currently available.

For wild fish extraction ($W$), the sustainable long-term target is
unlikely to be zero. Atlantic salmon require long-chain omega-3 fatty
acids and high-quality protein sources that still depend partially on
marine-derived ingredients. Although the marine ingredient share of feed
has fallen substantially, reaching near-zero $W$ would likely require
large-scale adoption of algae-derived omega-3 oils, insect proteins,
or other novel feed technologies that are not yet available at sufficient
scale or cost to fully replace marine inputs. The methods currently
driving decoupling, feed efficiency improvements, substitution toward
plant ingredients, and greater use of fish-processing by-products,
can probably reduce $W$ further, but are unlikely to eliminate it
entirely under current production systems.

For persistent antiparasitic discharge ($A_{out}$), a near-zero level
appears more realistic. The sharp decline observed after 2015 shows
that non-chemical treatments and oral treatments can dramatically reduce
persistent discharge. However, maintaining near-zero discharge in the
long term depends on whether sea lice can be controlled without renewed
dependence on persistent bath treatments. The main technological
constraint is the absence of a commercially viable and highly effective
sea lice vaccine. The methods currently used and recommended in this
report, preventive barriers, mechanical removal systems, cleaner fish,
and vaccine research, could plausibly reduce persistent discharge to
near-zero levels, but this outcome is not yet guaranteed.


\paragraph{Strengthen and extend the national monitoring infrastructure:}
Both decouplings documented in this report were made possible in part
by the existence of granular, publicly accessible national data on feed
composition, lice counts, treatment events, and substance quantities.
Without the Nofima benchmark series, the BarentsWatch \citep{barentswatch} lice reporting
system, and the NVI Fish Health Reports \citep{nvi2019}, neither IPAT decomposition
could have been constructed. This infrastructure should not be taken for
granted. We recommend that Norwegian authorities commit to continuing
and extending these data collection efforts: in particular by publishing
updated Nofima feed benchmarks beyond the 2013 endpoint of our analysis,
and by maintaining consistent treatment classification categories in the
NVI reports so that future researchers can extend IPAT~2 without the
reclassification discontinuities encountered around the 2017 to 2018 cleaner
fish data. Transparent and consistent monitoring is not a byproduct of good
environmental governance; in the Norwegian case, it has been a precondition
for it.




%==================================================
% 10. CONCLUSION
%==================================================
\chapter{Conclusion}

This report set out to answer two questions about the environmental footprint of Norwegian salmon aquaculture: whether two biodiversity pressures declined despite production growth, what factors explain the trends, and what role the regulatory framework played. Our two IPAT decompositions provide clear and 
quantitative answers to all three: absolute decoupling has occurred on both indicators, and by a large margin.

For wild fish extraction, $W$ fell by 64\% between 2000 and 2013 while production more than doubled. The decoupling was driven by three simultaneous efficiency improvements in $F/P$, $M/F$, and $W/M$, which together more than offset the effect of production growth. For persistent antiparasitic discharge, $A_{out}$ fell by 95\% between 2012 and 2019 while production continued to grow. This was dominated by the structural shift from bath treatments to oral and non-chemical alternatives, captured by $\delta$ and $\phi$.

The two decouplings were shaped by different forces. The decline in $W$ was market-led: economic pressure to reduce the cost of marine commodities, combined with certification standards and feed industry R\&D, drove improvements in $M/F$ and $W/M$ without any regulatory mandate on feed composition. The decline in $A_{out}$ required a more active regulatory hand such as the mandatory lice thresholds and the Traffic Light System were essential in producing the near-complete elimination of persistent chemical discharge. Market forces alone would likely not have produced a transition of this speed and magnitude.

Three caveats apply. First, each IPAT captures one pressure in isolation, a reduction in $W$ says nothing about the land-use implications of soy-based substitutes, and a reduction in $A_{out}$ says nothing about the welfare costs of mechanical delousing. Second, IPAT~1 relies on linear interpolation between four benchmark points, and IPAT~2 uses national aggregate data that mask regional variation. Third, both study periods end before 2020, and it is not guaranteed that the efficiency trends documented here have been maintained since.

Norway's experience offers a rare case study in absolute decoupling between 
industrial growth and environmental pressure. It required market incentives, regulatory pressure, technological innovation, and a data infrastructure that made outcomes visible and accountable. Identifying which conditions are load-bearing, as our decomposition framework helps to do, is a prerequisite for reproducing comparable outcomes elsewhere.

\newpage
%==================================================
% APPENDIX
%==================================================
\appendix

\chapter{Data and Code}

\section{Data Sources Summary}

Table~\ref{tab:datasources} summarises the primary datasets used in 
both IPAT decompositions, including their source, period covered, 
and the variables they provide.

\begin{table}[H]
\centering
\caption{Primary datasets used in the IPAT decompositions}
\begin{tabular}{llll}
\hline
\textbf{Dataset} & \textbf{Variables} & \textbf{Period} & \textbf{Source} \\
\hline
Fiskeridir sales statistics & $P$ (tonnes) & 2000--2019 & Fiskeridir (2024) \\[4pt]
Nofima benchmark reports & $F/P$, $M/F$, $W/M$ & 2000--2013 & Ytrestøyl et al. (2015) \\[4pt]
BarentsWatch Fish Health API & $N_{lice}/N_{salmon}$, $Tr_{total}$, $Tr_{chem}$ & 2012--2019 & BarentsWatch (2024) \\[4pt]
Jensen et al. (2020) & $A_{dose}$ by substance & 2012--2019 & Myhre Jensen et al. (2020) \\
\hline
\end{tabular}
\label{tab:datasources}
\end{table}

\section{IPAT 1 Data: Wild Fish Extraction}

Table~\ref{tab:ipat1data} reports the annual values of all four 
factors in the IPAT~1 decomposition, together with the computed 
wild fish extraction $W$.

\begin{table}[H]
\centering
\caption{IPAT 1 annual data, 2000--2013}
\begin{tabular}{crrrrr}
\hline
\textbf{Year} & \textbf{$P$ (t)} & \textbf{$F/P$} & \textbf{$M/F$} & \textbf{$W/M$} & \textbf{$W$ (t)} \\
\hline
2000 & 439,874 & 1.50 & 0.646 & 1.80 & 767,228 \\
2010 & 1,008,147 & 1.20 & 0.414 & 1.10 & 551,501 \\
2012 & 1,232,094 & 1.15 & 0.307 & 0.70 & 304,634 \\
2013 & 1,168,324 & 1.15 & 0.292 & 0.70 & 274,626 \\
\hline
\end{tabular}
\label{tab:ipat1data}
\end{table}

\noindent Intermediate years (2001--2009, 2011) are obtained by 
piecewise linear interpolation between benchmark points, as described 
in Section~4.

\section{IPAT 2 Data: Antiparasitic Chemical Discharge}

Table~\ref{tab:ipat2data} reports the annual values of all five 
factors in the IPAT~2 decomposition, together with the computed 
discharge $A_{out}$.

\begin{table}[H]
\centering
\caption{IPAT 2 annual data, 2012--2019.
The identity is $A_{out} = N_{\text{salmon}} \times \lambda \times \tau \times \phi \times \delta$.}
\begin{tabular}{crrrrrrr}
\hline
\textbf{Year}
  & \textbf{$N_{\text{salmon}}$}
  & \textbf{$\lambda$}
  & \textbf{$\tau$ ($\times 10^{-5}$)}
  & \textbf{$\phi$}
  & \textbf{$\delta$}
  & \textbf{$A_{out}$ (kg)} \\
  & (fish)
  & (lice/fish)
  & (events/lice)
  & (dim'less)
  & (kg/event)
  & \\
\hline
2012 & 246\,418\,984 & 0.214 &  9.91 & 0.500 & 1.719 & 4\,484 \\
2013 & 233\,664\,724 & 0.146 & 14.52 & 0.502 & 1.421 & 3\,535 \\
2014 & 251\,671\,172 & 0.207 & 10.11 & 0.460 & 2.147 & 5\,186 \\
2015 & 260\,669\,156 & 0.167 & 14.79 & 0.380 & 1.791 & 4\,364 \\
2016 & 246\,723\,848 & 0.163 & 16.30 & 0.356 & 0.701 & 1\,638 \\
2017 & 247\,270\,552 & 0.146 & 19.65 & 0.153 & 0.269 &   293 \\
2018 & 256\,400\,642 & 0.142 &  8.38 & 0.239 & 0.259 &   188 \\
2019 & 272\,808\,408 & 0.159 &  7.98 & 0.274 & 0.223 &   212 \\
\hline
\end{tabular}
\label{tab:ipat2data}
\end{table}

\section{Code}

All data processing and figures were produced in python and R.
The full source code and datasets are publicly available at
\url{https://github.com/htthomasepfl/hum470-aquaculture-biodiversity}. 
The key computational steps are summarised below.

\subsection{IPAT 1 — Wild Fish Extraction (R)}

\begin{verbatim}
# Nofima benchmark values
years_bm <- c(2000, 2010, 2012, 2013)
FCR  <- c(1.50, 1.20, 1.15, 1.15)
MF   <- c(0.646, 0.414, 0.307, 0.292)
WM   <- c(1.80, 1.10, 0.70, 0.70)

# Piecewise linear interpolation
years <- 2000:2013
FCR_i <- approx(years_bm, FCR, xout = years)$y
MF_i  <- approx(years_bm, MF,  xout = years)$y
WM_i  <- approx(years_bm, WM,  xout = years)$y

# Wild fish extraction
W <- P * FCR_i * MF_i * WM_i
\end{verbatim}

\subsection{IPAT 2 — Antiparasitic Discharge (Python)}

\begin{verbatim}
# Discharge calculation (H2O2 excluded)
bath_kg = (A_azamethiphos + A_cypermethrin 
           + A_deltamethrin)
oral_kg = (A_diflubenzuron + A_teflubenzuron 
           + A_emamectin)
A_out   = bath_kg * 1.00 + oral_kg * 0.03

# Ln-additive decomposition
import numpy as np
ln_contrib = np.diff(np.log([x2012, x2019]))
\end{verbatim}







%==================================================
% BIBLIOGRAPHY
%==================================================
%\newpage
\bibliographystyle{apalike}
\bibliography{bibliography}


% \begin{thebibliography}{99}

% \bibitem{sciencedirect2020b}
% Aaen, S. M., Helgesen, K. O., Bakke, M. J., Kaur, K., and Horsberg, T. E. (2015).
% Drug resistance in sea lice: a threat to salmonid aquaculture.
% \textit{Trends in Parasitology}, 31(2):72--81.
% \url{https://doi.org/10.1016/j.pt.2014.12.006}

% \bibitem{barrett2020}
% Barrett, L. T., Overton, K., Stien, L. H., Oppedal, F., and Dempster, T. (2020).
% Effect of cleaner fish on sea lice in Norwegian salmon aquaculture: a national scale data analysis.
% \textit{International Journal for Parasitology}, 50(4):277--283.
% \url{https://doi.org/10.1016/j.ijpara.2019.12.005}

% \bibitem{coolchoices2024}
% Cohen-Dumani, J., Billard, E., Alhaj Kaddour, Z., and Solomon, K. (2024).
% Cool choices in a heating world: The challenges of the cooling industry, from ozone depleting to climate warming refrigerants.
% Technical report, EPFL, SHS HUM-471 Economic Growth and Sustainability.
% Course report, Prof. Philippe Thalmann.

% \bibitem{fni2023}
% Fridtjof Nansen Institute (2023).
% Norway has the strictest aquaculture regulations in the world. Will other countries follow suit?
% \textit{Science Norway}.
% \url{https://partner.sciencenorway.no/aquaculture-fish-farm-fish-farming/norway-has-the-strictest-aquaculture-regulations-in-the-world-will-other-countries-follow-suit/2095202}

% \bibitem{myhrjensen2020}
% Myhre Jensen, P., Horsberg, T. E., Samuelsen, O. B., and Aaen, S. M. (2020).
% Anti-sea lice agents in Norwegian aquaculture; surveillance, treatment trends and possible implications for food safety.
% \textit{Aquaculture}, 519:734385.
% \url{https://doi.org/10.1016/j.aquaculture.2019.734385}

% \bibitem{nhm2024}
% Natural History Museum (2024).
% The problem of sea lice in salmon farms.
% Natural History Museum, London.
% \url{https://www.nhm.ac.uk/discover/the-problem-of-sea-lice-in-salmon-farms.html}

% \bibitem{nes2025}
% Nes, S. et al. (2025).
% Salmon lice biology, environmental factors, and smolt behaviour with implications for the Norwegian salmon farming management system: A critical review.
% \textit{Reviews in Aquaculture}.
% \url{https://doi.org/10.1111/raq.12953}

% \bibitem{barentswatch}
% Norwegian Coastal Administration and Partners (2024).
% BarentsWatch: Norwegian Coastal Monitoring Portal.
% \url{https://www.barentswatch.no/data}

% \bibitem{fiskeridir2024}
% Norwegian Directorate of Fisheries (2024).
% Statistics for Aquaculture: Atlantic Salmon and Rainbow Trout.
% Norwegian Directorate of Fisheries (\textit{Fiskeridirektoratet}).
% \url{https://www.fiskeridir.no/english/aquaculture/statistics-for-aquaculture/atlantic-salmon-and-rainbow-trout}

% \bibitem{nvi2019}
% Norwegian Veterinary Institute (2019).
% Fish Health Report 2019.
% Norwegian Veterinary Institute (\textit{Veterinærinstituttet}).
% \url{https://www.vetinst.no/rapporter-og-publikasjoner/rapporter/2020/fish-health-report-2019}

% \bibitem{oceanconservation}
% Ocean Conservation Science (2024).
% Forage Fish.
% \url{https://www.oceanconservationscience.org/forage-fish}

% \bibitem{samuelsen2020}
% Samuelsen, O. B., Parsons, A. E., Husa, V., Hannisdal, R., Tjensvoll, T., Reinardy, H. C., and Evenset, A. (2020).
% The impact of anti-sea lice pesticides, azamethiphos and deltamethrin, on European lobster (\textit{Homarus gammarus}) larvae in the Norwegian marine environment.
% \textit{Environmental Pollution}, 262:114509.
% \url{https://doi.org/10.1016/j.envpol.2020.114509}

% \bibitem{plos2018}
% Sandvik, A. D. et al. (2018).
% Evaluation of a national operational salmon lice monitoring system --- From physics to fish.
% \textit{PLOS ONE}, 13(7):e0201338.
% \url{https://doi.org/10.1371/journal.pone.0201338}

% \bibitem{seafoodsource2024}
% SeafoodSource (2024).
% How does Norway's traffic light system for salmon farming work?
% \textit{SalmonBusiness}.
% \url{https://www.salmonbusiness.com/how-does-norways-traffic-light-system-for-salmon-farming-work/}

% \bibitem{sunstone2026}
% Sunstone Institute (2026).
% Aquaculture Data Report: Eutrophication Pressure from Norwegian Fish Farming.
% \url{https://sunstone.institute/datapackages/aquaculture-data}

% \bibitem{frontiers2024}
% Sverdrup-Jensen, S. et al. (2024).
% Conceptualization of the Norwegian feed system of farmed Atlantic salmon.
% \textit{Frontiers in Marine Science}, 11:1378970.
% \url{https://doi.org/10.3389/fmars.2024.1378970}

% \bibitem{tacon2008}
% Tacon, A. G. J. and Metian, M. (2008).
% Global overview on the use of fish meal and fish oil in industrially compounded aquafeeds: Trends and future prospects.
% \textit{Aquaculture}, 285(1--4):146--158.
% \url{https://doi.org/10.1016/j.aquaculture.2008.08.015}

% \bibitem{fishsite2024}
% The Fish Site (2024).
% The innovative solutions being used to tackle sea lice on salmon farms.
% \url{https://thefishsite.com/articles/the-innovative-solutions-being-used-to-tackle-sea-lice-on-salmon-farms}

% \bibitem{guardian2026}
% The Guardian (2026).
% Norwegian fish farms polluting fjords with waste likened to raw sewage of millions of people.
% \url{https://www.theguardian.com/world/2026/may/04/norwegian-fish-farms-polluting-fjords-with-waste-likened-to-raw-sewage-of-millions-of-people}

% \bibitem{nofima2014}
% Ytrestøyl, T., Aas, T. S., and Åsgård, T. (2014).
% Utilisation of feed resources in production of Atlantic salmon in Norway.
% Technical Report 36/2014, Nofima.

% \bibitem{ytrestoeyl2015}
% Ytrestøyl, T., Aas, T. S., and Åsgård, T. (2015).
% Utilisation of feed resources in production of Atlantic salmon (\textit{Salmo salar}) in Norway.
% \textit{Aquaculture}, 448:365--374.
% \url{https://doi.org/10.1016/j.aquaculture.2015.06.023}

% \end{thebibliography}


\section*{AI Assistance Disclosure}

AI assistance (Claude, Anthropic) was used in the preparation of this report for the following purposes:

\begin{enumerate}
    \item Reformatting and copy-editing prose for clarity and concision.
    \item Reducing paragraph length to meet page constraints.
\end{enumerate}

AI was not used to generate research questions, select data sources, perform any quantitative analysis, produce any figures or tables, or formulate any interpretations or conclusions. All analytical and interpretive content is the authors' own work.

\end{document}